\section{Related Work}\label{sec:related}

\paragraph{Zero-knowledge proof security}
\zkap~\cite{zkap} applies static analysis to Circom circuits; Picus~\cite{picus} verifies Circom witness uniqueness via SMT; zkFuzz~\cite{zkfuzz} combines constraint-level coverage with symbolic witness generation for Circom. All three target the Circom ecosystem ($10^3$--$10^4$ constraints). zkVM arithmetizations differ in two ways: \emph{scale} ($10^5$--$10^6$ constraints with cross-linked trace tables and nonlinear arithmetic over 254-bit prime fields, beyond SMT solver capacity) and \emph{specification} (verifying Type~II soundness requires modeling both constraint system and RISC-V semantics---no production zkVM provides such a specification). \tool's differential approach sidesteps the specification requirement: semantic disagreement across backends constitutes evidence of a bug. FORAY~\cite{foray} synthesizes DeFi attacks at the application level; Tabby~\cite{tabby} optimizes circuit \emph{efficiency} via program synthesis, complementary to \tool's focus on \emph{correctness}.

\arguzz~\cite{arguzz} is the closest work: it applies metamorphic testing with fault injection to zkVMs. However, \arguzz lacks a semantic model of constraint-level boundary conditions and operates exclusively as a differential detector (analogous to \tool's Type~I detector). Our dual-detector characterization (\autoref{sec:formal-properties}) shows this is inherently limited to Type~I bugs. On the four backends where \arguzz has full support, \tool achieves \tbd{$1.5\times$} higher recall (\tbd{17 vs.\ 11}); \tool's detections are a strict superset. The additional bugs are predominantly Type~II, confirming that witness injection---not just better mutation---drives the improvement. \autoref{tbl:approach-comparison} summarizes key differences.

\begin{table}[t]
\caption{Comparison of zkVM vulnerability detection approaches.}
\label{tbl:approach-comparison}
\centering
\resizebox{\columnwidth}{!}{%
\footnotesize
\begin{tabular}{lccccc}
\toprule
\textbf{Dimension} & \textbf{\zkap} & \textbf{Picus} & \textbf{zkFuzz} & \textbf{\arguzz} & \textbf{\tool} \\
\midrule
Approach & Static & Symbolic & Hybrid & Dynamic & Dynamic \\
Target & Circom & Circom & Circom & zkVM & zkVM \\
Feedback & None & N/A & Constr.\ cov. & Metamorphic & Sig.\ novelty \\
ISA-aware & No & No & No & No & Yes \\
Witness inject & No & No & Symbolic & No & Semantic \\
Spec-free oracle & No & No & No & Yes & Yes \\
Cross-backend & No & No & No & Partial & 6 backends \\
\bottomrule
\end{tabular}%
}
\end{table}

\paragraph{Domain-specific coverage metrics}
\tool belongs to a growing trend replacing edge coverage with domain-specific metrics. Data Coverage~\cite{datacoverage} (USENIX Security 2024) replaces code coverage with data-flow coverage for database-system testing, observing that bugs cluster at specific data patterns rather than code paths. zkFuzz~\cite{zkfuzz} (IEEE S\&P 2026) uses constraint-level coverage for Circom testing. \tool extends this principle to RISC-V zkVMs with two key differences: (1)~categories serve a \emph{dual role} as both a coverage metric and a targeting mechanism (anchoring witness injection); and (2)~\tool's differential oracle does not require a formal specification, enabling it to scale across six heterogeneous backends without per-backend symbolic modeling.

\paragraph{Differential testing, fuzzing, and mutation scheduling}
Differential testing~\cite{mckeeman1998differential,csmith} compares outputs of multiple implementations; \tool adapts this to zkVMs with signature-novelty feedback. Coverage-guided fuzzers (AFL~\cite{afl}, LibFuzzer~\cite{libfuzzer}) use edge coverage, inadequate for zkVMs where interesting behaviors are at the constraint level. Grammar-aware fuzzers (Nautilus~\cite{nautilus}, Gramatron~\cite{gramatron}) ensure structural validity; \tool's ISA-aware mutations additionally incorporate semantic boundary awareness. MOpt~\cite{mopt} and EcoFuzz~\cite{ecofuzz} apply optimization-based scheduling to mutation operators; NEZHA~\cite{nezha} uses $\delta$-diversity for differential testing; SemFuzz~\cite{semfuzz} extracts semantic information from CVE descriptions. \tool's UCB1 bandit couples domain-specific category feedback with mutation selection, and uniquely provides a formal bug-space characterization (Type~I/II) where the same categories drive both the bandit's reward and the injector's targeting.
